\section{Balancing Innovation and Execution}

A \gls{CTO} must balance the need for new features with the long-term health of the codebase. This involves managing technical debt and making strategic build-versus-buy decisions.

\subsection{Managing Technical Debt}

Technical debt is an inevitable part of software development. You must manage it actively to avoid slowing down your team.

\begin{table}[h]
    \centering
    \begin{tabular}{|c|c|}
        \hline
        \textbf{Type of Technical Debt} & \textbf{Impact}                                   \\
        \hline
        Intentional                     & Short-term gains, long-term costs                 \\
        \hline
        Accidental                      & Results from poor decisions or lack of knowledge  \\
        \hline
        Environmental                   & Arises due to technology changes and system aging \\
        \hline
    \end{tabular}
    \caption{Types of Technical Debt}
\end{table}

Key strategies for managing technical debt include:

\begin{itemize}
    \item Regular refactoring and addressing debt incrementally.
    \item Prioritizing high-impact debt resolution.
    \item Educating stakeholders on long-term costs of accumulated debt.
\end{itemize}

\subsection{Evaluating Build vs. Buy Decisions}

Decide when to build custom solutions and when to leverage third-party tools. Consider:

\begin{itemize}
    \item \textbf{Cost:} Initial development versus long-term maintenance.
    \item \textbf{Time-to-Market:} Third-party tools can accelerate launch.
    \item \textbf{Customization Needs:} Does the existing solution meet your business requirements?
    \item \textbf{Vendor Lock-in:} What are the risks of depending on a third-party provider?
\end{itemize}
